% ============================================================
% RLHF on a Shoestring — report đồ án
% Biên dịch: pdfLaTeX (Overleaf: Menu -> Compiler -> pdfLaTeX)
% Nội dung đồng bộ với report.md ở repo root — sửa bên nào thì cập nhật bên kia.
% Khi trường có template riêng: bê phần \section vào template, giữ references.bib.
% ============================================================
\documentclass[12pt,a4paper]{article}

\usepackage[utf8]{inputenc}
\usepackage[T5]{fontenc}
\usepackage[vietnamese]{babel}
\usepackage[a4paper,margin=2.5cm]{geometry}
\usepackage{booktabs}        % bảng đẹp: \toprule \midrule \bottomrule
\usepackage{tabularx}        % bảng tự co giãn chiều rộng
\usepackage{graphicx}
\usepackage{enumitem}
\usepackage[colorlinks=true,linkcolor=blue,citecolor=blue,urlcolor=blue]{hyperref}
\usepackage[numbers,sort&compress]{natbib}
\usepackage{xcolor}

% Đánh dấu chỗ chờ số liệu thật — tìm \tbd trong file để biết còn thiếu gì
\newcommand{\tbd}[1]{\textcolor{red}{[TBD: #1]}}

\title{\textbf{RLHF on a Shoestring:\\ Building a Vietnamese Chatbot from a Raw Base Model on Free-Tier GPUs}}
\author{Đinh Phúc\thanks{\tbd{tên đầy đủ, MSSV, khoa/trường, GVHD}}}
\date{\today}

\begin{document}
\maketitle

\begin{abstract}
Đồ án xây dựng và vận hành trọn vẹn pipeline RLHF cho chatbot tiếng Việt, xuất phát từ một base model thô chưa qua instruction tuning (\texttt{Qwen/Qwen3-1.7B-Base}): Continual Pretraining (CPT) trên $\sim$400M token tiếng Việt, Supervised Fine-Tuning trên 12.7k hội thoại đa lượt, Reward Model, Direct Preference Optimization, đánh giá định lượng, và triển khai GGUF/Ollama. Toàn bộ huấn luyện chạy trên hạ tầng miễn phí (Kaggle T4 16GB, quota 30h/tuần) nhờ cơ chế checkpoint-resume qua HuggingFace Hub cho phép chia nhỏ mọi bước theo phiên. Đồ án trả lời bằng số liệu ba câu hỏi thiết kế: giá trị của CPT đối với tiếng Việt, mức độ thay thế của DPO đối với PPO, và hiệu quả của English replay trong việc chống catastrophic forgetting. \tbd{1--2 câu kết quả chính sau khi có số liệu từ E1--E3}
\end{abstract}

\section{Đặt vấn đề và động cơ}

Đa số đồ án ``chatbot tiếng Việt'' cùng cấp đi theo lối: lấy một model đã instruction-tuned sẵn (Llama-Instruct, Qwen-Instruct), fine-tune nhẹ, rồi demo. Cách đó ra sản phẩm nhanh nhưng bỏ qua toàn bộ phần khó và đáng học nhất: \emph{model biết chat từ đâu ra?} Đồ án này chọn con đường ngược lại --- bắt đầu từ raw weights chỉ biết dự đoán token kế tiếp, tự tay xây từng tầng biến nó thành trợ lý tiếng Việt, đúng quy trình công nghiệp mà InstructGPT~\citep{ouyang2022instructgpt} định hình và Sailor2~\citep{dou2025sailor2} áp dụng cho các ngôn ngữ Đông Nam Á.

Ràng buộc tự đặt: \textbf{chỉ dùng compute miễn phí} (Kaggle T4 16GB, phiên tối đa $\sim$9 giờ, quota 30 giờ/tuần). Ràng buộc này không phải điểm yếu phải giấu mà là một phần của đề bài: mọi kỹ thuật trong đồ án đều phải sống được trong điều kiện phiên chạy bị cắt bất kỳ lúc nào.

\section{Đóng góp}

Đóng góp của đồ án không nằm ở việc tạo ra một model vượt các LLM tiếng Việt hiện có --- điều bất khả thi ở quy mô tài nguyên sinh viên --- mà ở ba điểm:

\begin{enumerate}[leftmargin=*]
  \item \textbf{Pipeline hoàn chỉnh từ gốc.} Xây dựng và vận hành trọn vẹn chuỗi CPT $\to$ SFT $\to$ RM $\to$ DPO $\to$ Eval $\to$ Deploy từ base model thô, thay vì tinh chỉnh trên model đã instruction-tuned sẵn. Mỗi tầng có notebook độc lập, chạy được thật, kèm smoke test.
  \item \textbf{Ba câu hỏi thiết kế được trả lời bằng số liệu} (Mục~\ref{sec:experiments}) --- những câu các công bố lớn thường bỏ ngỏ ở quy mô nhỏ: (i) continual pretraining có đáng chi phí với tiếng Việt không; (ii) DPO thay được PPO đến đâu, với chi phí bao nhiêu; (iii) English replay 20\% có thực sự chặn được catastrophic forgetting không.
  \item \textbf{Quy trình khả tái lập trên hạ tầng miễn phí.} Cơ chế checkpoint-resume qua HuggingFace Hub (push nguyên checkpoint folder --- adapter, optimizer, scheduler, RNG --- mỗi \texttt{save\_steps}; tự dừng trước ngân sách giờ; phiên sau chạy lại là nối tiếp đúng step) biến mọi bước huấn luyện thành chuỗi phiên 8 giờ có thể lặp vô hạn. Sailor2 công bố cookbook cho người có hàng trăm GPU~\citep{dou2025sailor2}; đồ án này là cookbook cho người có một chiếc T4 mượn.
\end{enumerate}

Sailor2 được dùng làm \emph{mốc đối chiếu} trong đánh giá, không phải điểm xuất phát.

\section{Công trình liên quan}

\paragraph{Vai trò tương đối của pretraining và alignment.} LIMA~\citep{zhou2023lima} chỉ ra rằng phần lớn kiến thức và năng lực ngôn ngữ được học trong pretraining, còn alignment chủ yếu dạy \emph{format và style} (superficial alignment hypothesis); URIAL~\citep{lin2023urial} củng cố định lượng ở mức token. Đây là cơ sở lý thuyết cho việc giữ bước CPT thay vì trông cậy vào dữ liệu hội thoại ở các bước sau. Về kỹ thuật CPT, \citet{ibrahim2024simple} chứng minh LR re-warming/re-decaying kết hợp replay dữ liệu cũ đủ đạt chất lượng tương đương train lại từ đầu; \citet{gururangan2020dont} là công trình kinh điển về domain-adaptive pretraining; \citet{luo2023forgetting} đo đạc catastrophic forgetting và cho thấy model càng nhỏ quên càng nặng.

\paragraph{Công thức chuẩn cho ngôn ngữ bản địa.} Các model tiếng Việt và Đông Nam Á mạnh đều CPT trên corpus bản địa trước khi alignment, không có ngoại lệ: Sailor~\citep{dou2024sailor} và Sailor2~\citep{dou2025sailor2} (CPT 400B token SEA + 100B replay từ Qwen2.5), SeaLLMs~\citep{nguyen2023seallms}, VinaLLaMA~\citep{nguyen2023vinallama} (CPT 800B token trên Llama-2), PhoGPT~\citep{nguyen2023phogpt} (pretrain từ đầu 102B token). Base model của đồ án là Qwen3~\citep{qwen3report}.

\paragraph{RLHF và các phương án thay thế.} Chuỗi SFT $\to$ RM $\to$ PPO do \citet{christiano2017preferences,stiennon2020summarize,ouyang2022instructgpt} định hình, với phân tích thực nghiệm sâu về reward model trong \citet{bai2022hh}; PPO~\citep{schulman2017ppo} là thuật toán RL nền; OpenRLHF~\citep{hu2024openrlhf} là framework dùng cho bước RM và ablation PPO. DPO~\citep{rafailov2023dpo} rút gọn RLHF về một classification loss chỉ cần policy và reference --- lý do nó là con đường chính của đồ án trên single-GPU.

\paragraph{Điều kiện tồn tại trên T4.} LoRA~\citep{hu2021lora}, QLoRA~\citep{dettmers2023qlora} (NF4 4-bit + adapter fp16 + paged optimizer) và rsLoRA~\citep{kalajdzievski2023rslora} (scale $\alpha/\sqrt{r}$ cho rank cao) là bộ ba cho phép train 1.7B trên 16GB VRAM. Corpus tiếng Việt lấy từ FineWeb2~\citep{penedo2025fineweb2}, English replay từ FineWeb-Edu~\citep{penedo2024fineweb}.

\section{Pipeline}
\label{sec:pipeline}

Thiết kế chi tiết từng bước nằm trong spec của repo (Mục 6); notebook chạy được tại \texttt{notebooks/}. Bảng~\ref{tab:pipeline} tóm tắt.

\begin{table}[ht]
\centering\small
\caption{Pipeline 6 bước, tất cả xuất phát từ \texttt{Qwen3-1.7B-Base} thô.}
\label{tab:pipeline}
\begin{tabularx}{\textwidth}{@{}lXl@{}}
\toprule
\textbf{Bước} & \textbf{Nội dung} & \textbf{Output (HF Hub)} \\
\midrule
1. CPT & QLoRA $r{=}64$, $\sim$400M token (Wikipedia-vi + FineWeb2 \texttt{vie\_Latn} + 20\% FineWeb-Edu replay), 6000 step & \texttt{Qwen3-1.7B-vi-cpt} \\
2. SFT & QLoRA $r{=}16$, 12.7k hội thoại đa lượt, ChatML, loss chỉ trên lượt assistant & \texttt{Qwen3-1.7B-vi-sft} \\
3. RM & OpenRLHF \texttt{train\_rm}, preference data Vi-Alpaca & \texttt{Qwen3-1.7B-vi-rm} \\
4. DPO & TRL \texttt{DPOTrainer} (chính); PPO trên Modal (ablation) & \texttt{Qwen3-1.7B-vi-dpo} \\
5. Eval & VMLU + LLM-judge win-rate + RM score trên held-out & --- \\
6. Deploy & GGUF $\to$ Ollama + Open WebUI & \texttt{...-dpo-gguf} \\
\bottomrule
\end{tabularx}
\end{table}

\paragraph{Vì sao Qwen3 chứ không phải Llama.} Tokenizer Qwen ($\sim$152k vocab, train đa ngôn ngữ) cho fertility tiếng Việt thấp gần bằng tiếng Anh, trong khi vocab 32k của Llama-2 băm vụn âm tiết có dấu --- và tokenizer là thứ CPT không sửa được nếu không mở rộng vocab như SeaLLMs~\citep{nguyen2023seallms}. Thứ hai, với ngân sách token cố định, xuất phát cao cộng liều bồi nhẹ đi xa hơn xuất phát thấp cộng liều khổng lồ: VinaLLaMA cần 800B token để đưa Llama-2 lên chuẩn tiếng Việt~\citep{nguyen2023vinallama}, gấp 2000 lần ngân sách của đồ án.

\section{Thiết kế thực nghiệm}
\label{sec:experiments}

Ba thí nghiệm, mỗi thí nghiệm thay đổi đúng một biến (Bảng~\ref{tab:ablations}). Quy tắc đánh giá: held-out split chưa từng dùng để train RM/DPO/PPO; LLM-judge chấm hai chiều (đảo vị trí câu trả lời) để khử position bias~\citep{zheng2023judge}; VMLU~\citep{vmlu2025} đo kiến thức tiếng Việt; Sailor2-1B làm mốc đối chiếu tham khảo.

\begin{table}[ht]
\centering\small
\caption{Ba câu hỏi ablation và thước đo tương ứng.}
\label{tab:ablations}
\begin{tabularx}{\textwidth}{@{}lXXl@{}}
\toprule
\# & \textbf{Câu hỏi} & \textbf{Điều kiện so sánh} & \textbf{Thước đo} \\
\midrule
E1 & CPT có đáng không? & SFT-từ-CPT vs SFT-từ-base-thô & VMLU, win-rate \\
E2 & DPO thay được PPO? & DPO (T4) vs PPO (Modal), cùng preference data & win-rate, RM score, GPU-hour \\
E3 & Replay chống forgetting? & \texttt{eval\_vi\_loss} / \texttt{eval\_en\_loss} suốt 6000 step CPT & loss curve \\
\bottomrule
\end{tabularx}
\end{table}

\section{Kết quả}
\label{sec:results}

\tbd{Điền dần khi các bước hoàn thành --- mỗi bước xong dán bảng số và một đoạn nhận xét. Nguồn số liệu: \texttt{trainer\_state.json} trong các checkpoint trên Hub (CPT/SFT), script eval của bước 5.}

\subsection{E1 --- Giá trị của CPT} \tbd{bảng VMLU + win-rate}
\subsection{E2 --- DPO vs PPO} \tbd{bảng chất lượng + GPU-hour}
\subsection{E3 --- Catastrophic forgetting} \tbd{hình loss curve vi/en từ log CPT}

\section{Thảo luận phản biện}

\begin{description}[leftmargin=0pt,style=nextline]
  \item[``Của em hơn gì những model đang có?''] Không cạnh tranh chất lượng với model train bằng trăm GPU --- đồ án tái tạo \emph{cách các model đó được tạo ra} ở quy mô sinh viên, và đo đạc từng tầng. Phần ``hơn'' nằm ở Bảng~\ref{tab:ablations}: những con số mà việc dùng model có sẵn không bao giờ trả lời được.
  \item[``Sao không dùng luôn Sailor2?''] Dùng model sẵn thì không còn tầng nào để đo --- không trả lời được câu hỏi nghiên cứu nào. Sailor2 vẫn có mặt trong đồ án, đúng vai của nó: mốc đối chiếu ở phần eval.
  \item[``Model thua ChatGPT thì làm để làm gì?''] Đúng, và mục tiêu chưa bao giờ là thắng --- mục tiêu là kiểm soát từng tầng của quy trình. ChatGPT không cho xem đường loss của reward model.
  \item[``Sao chọn 1.7B nhỏ thế?''] Câu hỏi của đồ án là câu hỏi \emph{phương pháp}, không phải \emph{quy mô} --- mọi kết luận ablation đo được ở 1.7B. Scale lên là vấn đề tiền, không phải vấn đề khoa học.
  \item[``CPT có ích thật không hay làm theo trend?''] Không khẳng định chay --- E1 được thiết kế để trả lời chính câu này bằng số. Nếu delta nhỏ, đó cũng là một phát hiện trung thực được ghi nhận.
  \item[``Chạy Kaggle free thì nghiêm túc được không?''] Ràng buộc compute là một phần đề bài: mọi bước train đều checkpoint-resume qua Hub, tái lập 100\% từ notebook công khai. Khả tái lập là thứ nhiều công bố lớn còn thiếu.
\end{description}

\section{Reproducibility statement}

Toàn bộ notebook chạy trực tiếp trên Kaggle T4 free tier; dataset và checkpoint trung gian nằm trên HuggingFace Hub theo quy ước \texttt{<user>/Qwen3-1.7B-vi-\{cpt,sft,rm,dpo\}}; mọi seed cố định (42); log huấn luyện đầy đủ trong \texttt{trainer\_state.json} của từng checkpoint. Một người khác với tài khoản Kaggle và HF token có thể tái tạo từng bước bằng Run All.

\bibliographystyle{plainnat}
\bibliography{references}

\end{document}
